% !TeX root = ../thesis.tex
\section*{Abstract}

Final accuracy reports only whether a Text-to-SQL system produced 
the expected result and does not describe 
the SQL attempts and errors preceding its answer.
Building on previous work, this study examined how recorded SQL attempts
and database feedback reveal recovery, repeated failures and differences
between models answering the same questions.

An evaluation methodology and supporting software were developed to analyse
these records alongside saved query results.
The selected models (Amazon Nova Lite and Claude Haiku 4.5) each completed 
two runs of a custom Chinook benchmark through a cloud-based
agentic system.

The main analysis covered 160 standalone questions per run,
which provided the primary quantitative evidence.
Query results were compared with reference results according to a defined set of criteria.
SQL trajectories were classified as direct successes,
successful recoveries from errors or unsuccessful outcomes,
with errors and subsequent modifications also analysed.
Supplementary clarification, two-turn interaction and safety cases were assessed
separately.

Haiku achieved an average final accuracy of 91.56\%, and Nova achieved 78.44\%.
Both models recovered from execution errors on five standalone questions per run.
Haiku recovered on every standalone question where it encountered an execution error.
Nova, however, left most such cases unresolved.
It also produced longer failed sequences, including SQL revisions and repeated errors.
In some cases, models reached correct results on the same questions
through different paths, such as direct success versus recovery.
Supplementary cases showed that unchanged interactive accuracy
could conceal additional recovery, while equal no-query rates concealed differences
in manually assessed safety responses.

This study contributes an evaluation method and software for examining system performance.
The use of the recorded trajectories allows for moving beyond final accuracy analysis
and may enhance debugging and model assessment for specific deployments.

\textbf{Keywords:} Text-to-SQL, Large Language Models, Agentic Systems,
Trajectory Evaluation, SQL Error Recovery, Execution Traces.